% TOP_PROBLEM: {"id": 3800019, "title": "Generating random simple polygons", "category_id": 6, "category": {"id": 6, "name": "geometry", "display_name": "Geometry", "description": "Euclidean and non-Euclidean geometry, geometric structures.", "slug": "geometry", "order_index": 6, "created_at": "2026-07-31T15:26:25.671Z"}, "status": "partially_solved", "problem_number": "AMR-037-0019", "set_id": 13}
% FIRST_POSTED: 2026-10-01
% ENTRY_KIND: statement_only
% UnsolvedMath ID 3800019; source catalogue code AMR-037-0019
% !TeX program = lualatex
% Definition and sources only; this file is not an LLM solution attempt.
\documentclass[11pt,a4paper]{article}
\usepackage[margin=25mm]{geometry}
\usepackage{fontspec}
\setmainfont{lmroman10-regular.otf}[BoldFont=lmroman10-bold.otf,
 ItalicFont=lmroman10-italic.otf,BoldItalicFont=lmroman10-bolditalic.otf]
\usepackage{amsmath,amssymb,mathtools}
\usepackage{unicode-math}
\setmathfont{latinmodern-math.otf}
\AtBeginDocument{\let\setminus\smallsetminus}
\usepackage{xurl,hyperref}
\hypersetup{colorlinks=true,urlcolor=blue}
\setcounter{secnumdepth}{0}
\setlength{\parindent}{0pt}
\setlength{\parskip}{5pt}
\setlength{\emergencystretch}{3em}
\begin{document}
\section*{Generating random simple polygons}\label{3800019}
{\small UnsolvedMath ID 3800019 · AMR-037-0019 · Geometry · AMR Open Problem Lists}
\subsection{Definitions and mathematical statement}
Let $S=\{p_1,\ldots,p_n\}\subset\mathbb R^2$ be in general position
(no three points collinear). A polygonization of $S$ is a cyclic ordering
of all its points whose straight edges form a simple closed polygon.
Identify cyclic shifts and reversal, and write $N(S)$ for the number of
resulting undirected polygonizations.

The sampling question asks for a randomized algorithm which outputs each
polygonization with probability exactly $1/N(S)$ and whose expected
running time is polynomial in the input size. For a bit-complexity
formulation, the point coordinates are rational and their binary lengths
are included in that size. Exact uniformity is stronger than a distribution
which is merely close to uniform. A separate question is the computational
complexity of evaluating the integer $N(S)$ exactly.

The extremal part asks for the maximum of $N(S)$ over all $n$-point sets,
including its asymptotic growth. Corresponding extremal counting questions
replace polygonizations by triangulations of $\operatorname{conv}(S)$ using
every point, or by noncrossing straight-line Hamiltonian paths or spanning
trees on $S$. These are distinct families; their edge sets are counted
once, without extra choices of a starting vertex or orientation.
\subsection{Short English statement}
Can all simple polygons through a given point set be counted or sampled
uniformly in polynomial time? How large can these families, and related
families of noncrossing graphs, become?
\subsection{Sources}
\begin{itemize}
\item[S1] ULAM.AI and the UnsolvedMath Contributors, supplied problems.json, record 3800019 (AMR-037-0019), AMR Open Problem Lists. Catalogue identity and source transcription; CC BY 4.0.
\url{https://huggingface.co/datasets/ulamai/UnsolvedMath}.
\item[S2] Erickson - Open Problems in Computational Geometry, Generating random simple polygons. Original source indexed by AMR-037-0019.
\url{https://jeffe.cs.illinois.edu/open/}.
\item[S3] Jeff Erickson, Generating random simple polygons.
\url{https://jeffe.cs.illinois.edu/open/randompoly.html}.
\end{itemize}
\end{document}
